\section{Discussion}\label{sec:disc}

The benchmarks of Sections~\ref{sec:verif} and~\ref{sec:cv7} put the transition-map chart and the
ambient level set on the same problems, the same meshes, and the same tolerances, and report the
measured difference between them. We now read those results as a single statement: \emph{when} does a
boundary chart improve on a neural signed-distance function, and when does it not? Three subsections
answer that question in order. We first collect the verdict per benchmark in a capability matrix
(Section~\ref{sec:disc:matrix}). We then state the preconditions and trade-offs that bound the method
(Section~\ref{sec:disc:precond}). We close with the limitations the data expose, each stated together
with what it does and does not threaten (Section~\ref{sec:disc:limits}).

\subsection{A capability matrix: where the transition map is exercised}\label{sec:disc:matrix}

The closed-form suite CV-1\,--\,CV-4 was not built to flatter the chart. Each of those four problems is
convex, collinear, or superposable, so the boundary-to-boundary correspondence never comes into play.
That correspondence is the transition map $\tmap=\map{B}^{-1}\circ\map{A}$, which carries a point on
body $A$'s boundary chart $\map{A}$ to its mate on body $B$'s chart $\map{B}$. On a smooth convex pair
there is no nonconvex multi-arc contact to enumerate, no cusp to resolve, and no curved-surface
correspondence to invert, so the map reduces in effect to the identity and the chart gap collapses onto
the SDF gap. CV-1 (Hertz) presses a smooth convex indenter; CV-2 (Cattaneo--Mindlin) adds a tangential
load to that pre-formed Hertz patch; CV-3 (Brazilian) is a pure stress boundary-value problem with
prescribed pole tractions; CV-4 is a symmetry-reduced nine-disc unit cell with prescribed diametral
loads. On all four the radial inverse-chart gap equals the SDF gap, and the two representations are
interchangeable. We call CV-1\,--\,CV-4 the \emph{equivalence regime}: the neural geometry is verified
to reproduce the closed forms
\citep{johnson1985contact,timoshenko1970elasticity,mindlin1949compliance,hondros1959evaluation}, and the
chart neither beats nor loses to the level set. The discriminating evidence therefore lives in
CV-5\,--\,CV-7, where the geometry is cusped, fractal, or rough.

Table~\ref{tab:capability} collects the verdict for each benchmark. Its columns separate three
independent questions: what geometry the benchmark presents, whether the transition map is actually
exercised (rather than reducing to the SDF), and the measured chart-versus-level-set outcome.

% code: docs/contact_verification_manual.md::11.8 (capability matrix); benchmarks/contact/cv_numerical/cv7_transition_map_contact.py
\begin{table*}
\centering
\caption{Capability matrix for the verification and validation suite. CV-1\,--\,CV-4 form the
equivalence regime: the geometry is convex, collinear, or superposable, the transition map is not
exercised, and the radial chart and the neural SDF are equal. CV-5\,--\,CV-7 exercise the transition map
and are where the measured advantage lives. Gap errors are relative to the surface RMS where a percentage
is given. All numbers are measured, not targets.}
\label{tab:capability}
\small
\begin{tabular}{@{}llll@{}}
\toprule
Benchmark & Geometry & Transition map exercised? & Chart vs.\ level set (measured) \\
\midrule
CV-1 & smooth convex indenter & no (radial gap reduces to SDF) & equivalent; FEM Hertz $a(F)$--$E^{*}$ to $\sim$1.6\% \\
CV-2 & Hertz patch $+$ tangential load & no (friction layer on CV-1) & equivalent; $c/a$ to $\sim$5--11\% \\
CV-3 & disc, prescribed pole tractions & no (stress BVP, no detector) & n/a; centre stress 1.62\% / 0.58\% \\
CV-4 & nine-disc unit cell & no (prescribed diametral loads) & n/a; equibiaxial centre 0.15\% \\
\midrule
CV-5 & cusped non-convex superformula & \textbf{yes} (multi-arc, cusps) & chart gap $3.8{\times}10^{-3}L$ / $0.42^{\circ}$ vs.\ SDF $8{\times}10^{-3}L$ / degraded; dynamics 0.04\% \\
CV-6 & Koch-snowflake fractal & \textbf{yes} (recursive descent) & $O(1)$ storage vs.\ $O(9^{n})$ SDF; refinement ceiling measured \\
CV-7 & real rough rock joint & \textbf{yes} (height-chart detector) & active set 98.9\% vs.\ 95.8\%; strength $-35$ to $-61$\%; dilatancy $-98$\% \\
\bottomrule
\end{tabular}
\end{table*}

% code: docs/contact_verification_manual.md::11.8 (equivalence regime)
The equivalence-regime rows earn their place in the argument rather than padding it. The chart's
advantage is specific to hard geometry, and CV-1\,--\,CV-4 are not hard geometry. By reporting that the
chart buys nothing on a sphere, we make the wins on CV-5\,--\,CV-7 attributable to the geometric
representation and not to an incidental difference in the solver. Figure~\ref{fig:cv_summary} gathers the
whole suite: the equivalence-regime errors against the closed forms on one side, and the chart-over-SDF
separation on the cusped superformula on the other.

\subsection{Preconditions and trade-offs}\label{sec:disc:precond}

Two conditions bound where the transition-map detector applies, and one property bounds what it returns.
The detector is defined only on geometry that admits the chart it uses, and the gap it reports is sharp
in sign but enlarged in magnitude. Each boundary below is a geometric fact about the parametrization,
checked once, not a solver tolerance that one might tighten.

\subsubsection{Star-shapedness and single-valuedness}

% code: solvers/contact/chart_gap.py::evaluate_gap_chart; solvers/contact/profile_chart_2d.py::NeuralHeight1D
A radial chart represents a body as a single-valued radius along every direction, and that
representation holds only when each ray actually meets the boundary once. Formally, the radial chart
$\rho_X:\Sph^{d-1}\to\R^{+}$ sends a unit direction to the distance from the center $c_X$ to the
boundary; it is single-valued exactly where every ray from $c_X$ meets the boundary one time. A plane or
a half-space is not star-shaped about any finite center, so floors and rigid platens stay on the SDF
path, while the radial detector handles the star-shaped obstacles: the superformula cam, the discs, and
the spheres. The height chart $z=h_X(x)$ relaxes the need for a center, but trades it for a different
condition. It represents the surface as a graph, so the surface may carry no overhangs and no re-entrant
detail. A genuinely re-entrant surface needs the full decoder atlas with Newton inversion at chart
overlaps \eqref{eq:newton-invert}, which is the Stage-2 generalization and is not exercised here.

\subsubsection{The radial gap is conservative-large, sign-exact, and not eikonal}

% code: solvers/contact/chart_gap.py::evaluate_gap_chart; solvers/contact/supershape.py
The radial detector reads the gap \emph{along the ray from the center}, not perpendicular to the
surface, and that is the one property that bounds the value it returns. The two readings differ by the
cosine of the angle $\alpha$ between the ray direction $\dhat$ and the true surface normal at the radial
foot. The next equation states that relation and its immediate consequence: the radial gap is always at
least as large in magnitude as the perpendicular gap.
%
\begin{equation}
\gn^{\mathrm{rad}} = \frac{\gn^{\perp}}{\cos\alpha} + O(\gn^{2})
\quad\Longrightarrow\quad
\abs{\gn^{\mathrm{rad}}} \ge \abs{\gn^{\perp}}\ \text{always,}
\label{eq:disc:radbias}
\end{equation}
%
where $\gn^{\mathrm{rad}}$ is the radial gap, $\gn^{\perp}$ is the true normal-projected gap, and the
sign convention is $\gn<0\Leftrightarrow$ penetration throughout. The factor $1/\cos\alpha$ is the only
moving part, and three consequences of its being positive make the enlargement acceptable. First, the
\emph{sign is exact}: a positive multiplier cannot flip $\gn^{\perp}$, so the active set $\{\gn<0\}$ read
from the radial chart equals the active set of the true SDF on a star-shaped body. Second, the
\emph{magnitude is conservative-large}: an exterior point reads a marginally larger gap and a penetrating
point a marginally deeper penetration, so the penalty is at worst slightly stiffer, never slipped.
Third, $\gn^{\mathrm{rad}}$ is $C^{1}$ wherever $\rho_X$ is, but it is \emph{not} an eikonal distance
\citep{kolev2003level,liu2020implicit}: its gradient norm $\norm{\grad\gn}$ can reach $\sim 10^{4}$ in
deep concave valleys, where the radial normal rotates quickly. A closest-point refinement recovers the
perpendicular gap, but we use it only for verification, because the perpendicular foot jumps across the
medial axis and would leak energy if wired into the integrator. The height chart obeys the analogous
statement: the graph offset of \eqref{eq:height-gap} is the small-slope-exact normal-projected gap, so
its fidelity lives in slopes and normals, not in gap magnitude.

\subsubsection{Storage and query cost}

% code: docs/contact_verification_manual.md::8 (CV-6 cost scaling)
The fractal benchmark shows the storage trade-off most clearly. CV-6 mates two Koch snowflakes, each a
four-contraction iterated function system carried to recursion depth $n$. The chart \emph{is} the
generating rule, so its storage is $O(1)$, independent of $n$; a pruned recursive inside-test answers a
query in $O(\text{depth})$ work and plateaus at $\approx 21$ visited nodes for $n$ up to 12. A neural or
uniform level set has no such generating rule and must instead discretize the ambient field at the
finest scale it resolves, which costs $O(9^{n})$ uniform (or $O(4^{n})$ adaptive) storage to chase a
self-similar boundary (Figure~\ref{fig:koch_cost}). The level set does answer each query more cheaply,
through an $O(1)$ lookup, but its fixed capacity caps the geometry it can represent
\citep{rahaman2019spectral,sitzmann2020implicit}, and that ceiling is what CV-6 measures.

\subsubsection{The chart wins on the gap and the active set, not necessarily on tip normals}

% code: benchmarks/contact/cv_numerical/cv7_transition_map_contact.py
The most carefully measured trade-off comes from routing the height-chart detector into the FEM contact
loop on the real rough joint (CV-7, Phase~3). On $10^{5}$ query points the chart's active set agrees
with the analytic reference 98.9\% of the time, against 95.8\% for the ambient SDF, at $1.47\times$ the
per-query cost; the contact-driving gap RMSE is 4.2\% of the surface RMS for the chart against 44.7\% for
the SDF, a $10.5\times$ sharper quantity. The chart does not, however, win everywhere. Its random-Fourier
height gradient gives a \emph{noisier tip normal}: at the sharpest asperity peaks the SDF returns a
median tip-normal error of $0.8^{\circ}$ against the chart's $10.3^{\circ}$. This local loss does not
weaken the strength prediction, and the reason is what contact enforcement actually consumes. The
penalty force and the friction cone read the gap and the active set, not the pointwise normal at an
isolated tip \citep{wriggers2006computational,laursen2002computational}, so a locally noisier height-field
normal does not propagate into the response, whereas the smoothed \emph{gap} of the SDF does.

\subsection{Limitations}\label{sec:disc:limits}

Four limitations follow directly from the measurements. We state each one with what it does and does not
put at risk, so that the reader can see the boundary of the present scope without having to weigh a
defensive list.

\subsubsection{The two-block friction residual}

% code: benchmarks/contact/cv_numerical/rock_joint_two_block.py; docs/cv7_session_results.md
The frictionless and one-block CV-7 solves converge cleanly: the contact Newton residual reaches
$\sim 10^{-9}$ with a consistent tangent. Two mutually deformable blocks under Coulomb friction do not.
Across four mesh refinements the friction residual holds at 3.8\%--4.8\% and does \emph{not} shrink with
refinement, which identifies it as a model property rather than a discretization error: it is the
non-smoothness of moving-master, node-to-surface Coulomb contact between two deforming rough faces
\citep{zavarise1998segment,fischer2005frictionless}, not a mortar formulation. What the residual does
\emph{not} threaten is the headline result, because the emergent quantities are mesh-converged: peak
$\muapp \to 0.53\pm 2\%$ and total dilation $\to 0.063$\,mm over the refinement sequence. The
$1.5$--$6\%$ residual band thus perturbs the friction balance without corrupting the converged response.
The path to a tight residual is a mortar or augmented-Lagrangian / semismooth-Newton treatment of the
tangential complementarity \citep{alart1991mixed,simo1992augmented,puso2004mortar,chouly2013nitsche},
which is a constraint-enforcement upgrade and is orthogonal to the geometric representation this paper
argues for.

\subsubsection{A single band-limited realization in small-strain elasticity}

% code: solvers/fem/rough_block_decoder.py; benchmarks/contact/cv_numerical/cv7_roughness_sweep.py
The rough-geometry validation rests on one band-limited surface realization, and the constitutive model
is small-strain linear elasticity. The decoder is band-limited by construction, through a finite
Fourier-feature bank, so the resolved roughness has a spectral cutoff, and a single realization is
statistically spiky and should not be read as an ensemble strength. The roughness sweep
(Figure~\ref{fig:cv7-roughness}) is the response to this concern: it varies the surface RMS and the
spectral cutoff and reports the dilatancy law $\muapp$ versus roughness as a trend, with peak $\muapp$
rising $0.61\to 1.20$ as the RMS grows $0.022\to 0.066$\,mm and the decoder reconstruction held below
8\%. The robust claim is therefore the trend, not any single absolute value
\citep{barton1977shear,barton1990review}. Finite strain, plasticity at asperity tips, and asperity
damage lie outside the present model \citep{ball1977convexity,ciarlet1988mathematical}.

\subsubsection{The contact gap is a small-slope vertical proxy}

% code: solvers/contact/profile_chart_2d.py::vertical_gap; docs/contact_verification_manual.md::11.9
On the height-chart joint the contact gap is the vertical offset $g_{\mathrm{vert}}=z-h_\theta(x)$,
projected onto the small-slope normal of \eqref{eq:height-gap}; here $z$ is the query height and
$h_\theta(x)$ is the learned surface height beneath it. This proxy is exact in sign and accurate to
$O(\text{slope}^{2})$ in magnitude, which is why the fidelity story for CV-7 is told in slopes and
normals rather than in gap magnitude. The proxy does not blunt the central claim, and the Inada
head-to-head shows why \citep{inada2020digitalrocks}. The height chart reconstructs the surface to
2.3\,\textmu m against the ambient SDF's 107\,\textmu m, and the SDF smooths the mean asperity angle
$19.4^{\circ}\to 12.5^{\circ}$; through the Patton law $\muapp=\tan(\phi_b+i)$ \citep{patton1966multiple}
that smoothing under-predicts peak strength by 61\%, while the \emph{total} dilation, set by the
large-scale waviness the SDF \emph{does} capture, changes by only $0.5\%$. The chart wins on the
slopes that set strength, then, and not on the low-frequency gap that both representations already get
right. Figure~\ref{fig:cv7-atlas-sdf} shows the same separation on the synthetic decoder geometry:
frictionless dilatancy emerges on the resolved surface, with apparent $\muapp$ rising $0\to 0.17$, and is
under-predicted by 98\% on the SDF-smoothed geometry, with frictional strength under-predicted by 35\%.

\subsubsection{The production oracle still uses the SDF gradient; the chart detector is opt-in}

% code: solvers/contact/gap.py::evaluate_gap; solvers/contact/contact_manager.py::body_gap_normal
The transition-map detector composes with the existing force stack behind a single flag. Penalty,
augmented-Lagrangian, and regularized Coulomb friction all consume only the $(\gn,\nbf,\norm{\fbf_N})$
contract \citep{kikuchi1988contact}, namely the gap, the surface normal, and the normal-force magnitude,
so the chart and the SDF are interchangeable at the force layer. The limitation is one of deployment, not of
principle: the production MPM oracle still computes the gap from the SDF gradient, and the chart detector
drives CV-5 analytically and CV-7 numerically (Phase~3) as an opt-in path (\code{detector="chart"})
rather than the default everywhere. Replacing the SDF throughout the dynamic solver is the stated
direction, and the explicit ChartMPM cross-check supports it by reproducing the Coulomb floor
($\muapp\approx 0.34$ against a base $\mu=0.4$), but it is not yet the default. What makes the eventual
swap conservative is the energy-consistency requirement that the matched normal of
\eqref{eq:matched-normal} be the gradient of the gap. The cyclic ledger of
Figure~\ref{fig:supp_cyclic} confirms it: the per-cycle balance
$W_{\mathrm{ext}}/(W_{\mathrm{fric}}+\Delta U_{\mathrm{el}}+W_{\mathrm{pen}}+W_{\mathrm{stick}})$ closes
from $1.09$ to $1.01$ on the deformable joint, which is the evidence that chart-driven contact is
energetically closed despite the friction residual above.

% code: benchmarks/contact/cv_numerical/rock_joint_decoder_cyclic.py
Taken together, the matrix and these caveats define the operating envelope of the method. The
transition-map chart is the representation of choice on star-shaped, single-valued, cusped, fractal, or
rough geometry that a fixed-capacity level set smooths; it is exactly equal to the SDF on the convex and
superposable problems; and it carries two costs that are the price of dropping the ambient field, a
conservative-large radial gap and, for now, an opt-in dynamic path.
